\section{Simulation Environment Design}

The simulation environment was developed using the OpenAI Gymnasium framework, which provides a standardized interface for reinforcement learning research and enables reproducible experimentation. The energy management problem is formulated as a sequential decision-making task and modeled as a Markov Decision Process (MDP), defined by the tuple $(S, A, P, R, \gamma)$, where $S$ denotes the state space, $A$ the action space, $P$ the transition dynamics, $R$ the reward function, and $\gamma$ the discount factor.

The state space ($S$) is designed to capture the essential operational information required for real-time energy management. At each time step $t$, the observation vector includes the current battery state of charge ($SoC_t$), the instantaneous residential load demand ($P_{\text{load},t}$), the available photovoltaic generation ($P_{\text{pv},t}$), and the current Time-of-Use electricity price ($\text{Price}_t$). In addition, to enable the agent to learn periodic daily patterns and anticipate upcoming price and demand variations, cyclical time features are incorporated using the sine and cosine encoding of the current hour. All state variables are normalized to the range $[0,1]$ to improve numerical stability and facilitate efficient neural network training.

The action space ($A$) is defined as a continuous domain to allow fine-grained control of the battery power. At each time step, the agent outputs a single scalar action $a_t \in [-1,1]$, representing the normalized charging or discharging command. Positive values correspond to battery charging, while negative values correspond to battery discharging. The actual battery power set-point is obtained by scaling the normalized action by the maximum inverter power rating ($P_{\max}$). A continuous action formulation is preferred over a discrete one, as it avoids quantization effects and reduces abrupt switching behavior, which can lead to inefficient operation and accelerated battery degradation.

The transition dynamics ($P$) are governed by the physical behavior of the battery energy storage system. The battery state of charge evolves as a function of the applied power command, the battery energy capacity, the time-step duration, and the charging and discharging efficiencies. Operational constraints, including minimum and maximum allowable state-of-charge limits, are strictly enforced within the environment. If an agent’s action would result in a violation of these constraints, it is clipped to the nearest feasible value. This constraint-handling mechanism ensures safe operation and reflects the hard physical limits present in real-world battery systems.

While the environment assumes that the current state observation provides sufficient information for decision-making, it is acknowledged that future load and generation values are not directly observable. As such, the formulation represents a practical approximation of a Markov process, consistent with common practice in real-time energy management applications. The inclusion of time-encoding features partially mitigates this limitation by allowing the agent to learn recurring temporal patterns from data.
